\documentclass[12pt,a4paper]{report}

% ══════════════════════════════════════════════════════════════
% PACKAGES
% ══════════════════════════════════════════════════════════════
\usepackage[utf8]{inputenc}
\usepackage[T1]{fontenc}
\usepackage{lmodern}
\usepackage[french]{babel}
\usepackage{geometry}
\usepackage{graphicx}
\usepackage{amsmath,amssymb}
\usepackage{booktabs}
\usepackage{tabularx}
\usepackage{multirow}
\usepackage{longtable}
\usepackage{float}
\usepackage{caption}
\usepackage{subcaption}
\usepackage{listings}
\usepackage{xcolor}
\usepackage{hyperref}
\usepackage{fancyhdr}
\usepackage{titlesec}
\usepackage{tocloft}
\usepackage{setspace}
\usepackage{parskip}
\usepackage{enumitem}
\usepackage{tikz}
\usetikzlibrary{shapes.geometric, arrows.meta, positioning, fit, calc}

\graphicspath{{diagrams/}}

\geometry{
  top=2.5cm,
  bottom=2.5cm,
  left=3cm,
  right=2.5cm
}

\onehalfspacing

% ── Hyperref setup ──
\hypersetup{
  colorlinks=true,
  linkcolor=blue!60!black,
  citecolor=green!50!black,
  urlcolor=blue!70!black,
  bookmarks=true,
}

% ── Code listings style ──
\definecolor{codebg}{rgb}{0.95,0.95,0.97}
\definecolor{codegreen}{rgb}{0,0.5,0}
\definecolor{codegray}{rgb}{0.5,0.5,0.5}
\definecolor{codepurple}{rgb}{0.58,0,0.82}

\lstdefinestyle{mystyle}{
  backgroundcolor=\color{codebg},
  commentstyle=\color{codegreen},
  keywordstyle=\color{blue},
  numberstyle=\tiny\color{codegray},
  stringstyle=\color{codepurple},
  basicstyle=\ttfamily\footnotesize,
  breakatwhitespace=false,
  breaklines=true,
  captionpos=b,
  keepspaces=true,
  numbers=left,
  numbersep=5pt,
  showspaces=false,
  showstringspaces=false,
  showtabs=false,
  tabsize=2,
  frame=single,
  rulecolor=\color{gray!30},
}
\lstset{style=mystyle}

% ── Headers/Footers ──
\pagestyle{fancy}
\fancyhf{}
\setlength{\headheight}{40pt}
\fancyhead[L]{\leftmark}
\fancyhead[R]{\thepage}
\renewcommand{\headrulewidth}{0.4pt}

% ══════════════════════════════════════════════════════════════
% TITLE PAGE
% ══════════════════════════════════════════════════════════════
\begin{document}

\hypersetup{pageanchor=false}

\begin{titlepage}
\begin{center}

\vspace*{1cm}

{\Large \textbf{R\'{e}publique Alg\'{e}rienne D\'{e}mocratique et Populaire}}\\[0.3cm]
{\large Minist\`{e}re de l'Enseignement Sup\'{e}rieur et de la Recherche Scientifique}

\vspace{1.5cm}

\rule{\linewidth}{0.5pt}\\[0.5cm]
{\Huge \bfseries Conception et D\'{e}veloppement d'un CRM Intelligent}\\[0.3cm]
{\LARGE \bfseries Int\'{e}grant un Module de Pr\'{e}diction des Ventes}\\[0.2cm]
{\LARGE \bfseries Bas\'{e} sur le Machine Learning}\\[0.2cm]
{\Large \bfseries pour le E-Commerce}\\[0.3cm]
\rule{\linewidth}{0.5pt}

\vspace{1cm}

{\Large \textbf{Projet de Stage}}\\[0.3cm]

\vspace{1cm}

{\large \textbf{Entreprise d'accueil :} BAT PROJET ENGINEERING}\\[0.2cm]
{\large Repr\'{e}sent\'{e}e par Mr. ATALLAH Mohammed}

\vspace{1cm}

\begin{minipage}{0.45\textwidth}
\begin{flushleft}
{\large \textbf{R\'{e}alis\'{e} par :}}\\[0.3cm]
{\large Rezaiguia Soltane Tadj Eddine \textit{(DS)}}\\[0.15cm]
{\large Bellatreche Mohamed Amine \textit{(DS)}}\\[0.15cm]
{\large Khelifi Ayyoub \textit{(CS)}}\\[0.15cm]
{\large Brahim Soheib \textit{(AI)}}
\end{flushleft}
\end{minipage}
\hfill
\begin{minipage}{0.45\textwidth}
\begin{flushright}
{\large \textbf{Encadr\'{e} par :}}\\[0.3cm]
{\large Dr. Khellat Souad}
\end{flushright}
\end{minipage}

\vfill

{\large Ann\'{e}e Universitaire 2025--2026}

\end{center}
\end{titlepage}

\hypersetup{pageanchor=true}

% ══════════════════════════════════════════════════════════════
% FRONT MATTER
% ══════════════════════════════════════════════════════════════
\pagenumbering{roman}

\tableofcontents
\newpage

% ══════════════════════════════════════════════════════════════
% INTRODUCTION GENERALE
% ══════════════════════════════════════════════════════════════
\pagenumbering{arabic}

\chapter*{Introduction G\'{e}n\'{e}rale}
\addcontentsline{toc}{chapter}{Introduction G\'{e}n\'{e}rale}

Le commerce \'{e}lectronique conna\^{i}t une croissance exponentielle en Alg\'{e}rie, particuli\`{e}rement dans le mod\`{e}le \textit{Cash-on-Delivery} (COD), o\`{u} le paiement s'effectue \`{a} la livraison. Ce mod\`{e}le, bien qu'adapt\'{e} au contexte local o\`{u} les solutions de paiement \'{e}lectronique restent limit\'{e}es, engendre des d\'{e}fis majeurs : un taux d'\'{e}chec de livraison de 30 \`{a} 50\%, des retours co\^{u}teux qui \'{e}rodent les marges, et une gestion inefficace des commandes.

Face \`{a} ces probl\'{e}matiques, les syst\`{e}mes de gestion de la relation client (CRM) traditionnels s'av\`{e}rent insuffisants. Ils se limitent \`{a} la gestion passive des donn\'{e}es sans exploiter le potentiel pr\'{e}dictif qu'elles renferment. L'int\'{e}gration de l'intelligence artificielle, et plus sp\'{e}cifiquement du Machine Learning, dans ces syst\`{e}mes ouvre de nouvelles perspectives pour transformer les donn\'{e}es historiques en informations d\'{e}cisionnelles actionnables.

Le pr\'{e}sent travail s'inscrit dans cette d\'{e}marche et vise \`{a} concevoir et d\'{e}velopper un CRM intelligent pour le e-commerce COD, int\'{e}grant un module de pr\'{e}diction bas\'{e} sur le Machine Learning. Ce module comprend trois axes principaux :

\begin{enumerate}
  \item \textbf{La pr\'{e}diction du risque de livraison} : un ensemble de mod\`{e}les (CatBoost, LightGBM, XGBoost) pour \'{e}valuer la probabilit\'{e} de succ\`{e}s de chaque commande.
  \item \textbf{La segmentation client} : une approche hybride HDBSCAN (grands jeux de donn\'{e}es) / KMeans (petits jeux de donn\'{e}es) pour identifier automatiquement des groupes de clients aux comportements similaires.
  \item \textbf{La pr\'{e}vision de la demande} : un mod\`{e}le LightGBM avec covariables du calendrier islamique pour anticiper les tendances de vente futures.
\end{enumerate}

Le rapport suit une progression continue : pr\'{e}sentation de l'organisme d'accueil, cadrage du probl\`{e}me, choix m\'{e}thodologiques, conception/impl\'{e}mentation, validation IA, puis s\'{e}curit\'{e} et conclusion. Cette structure permet de conserver un fil conducteur commun entre enjeux m\'{e}tier, d\'{e}cisions techniques et r\'{e}sultats mesur\'{e}s.


% ══════════════════════════════════════════════════════════════
% CHAPITRE 1 : PRESENTATION DE L'ORGANISME D'ACCUEIL
% ══════════════════════════════════════════════════════════════
\chapter{Pr\'{e}sentation de l'Organisme d'Accueil}

\section{Profil de BAT PROJET ENGINEERING}

Le stage s'est d\'{e}roul\'{e} au sein de \textbf{BAT PROJET ENGINEERING}, structure partenaire du projet et environnement d'exp\'{e}rimentation du CRM. L'organisme intervient sur des activit\'{e}s de gestion op\'{e}rationnelle et de pilotage m\'{e}tier, avec un besoin concret d'outillage num\'{e}rique pour fiabiliser les processus.

Dans ce contexte, la probl\'{e}matique COD (Cash-on-Delivery) est apparue comme prioritaire : volume de commandes croissant, suivi manuel dispers\'{e}, difficult\'{e} \'{a} anticiper les risques d'\'{e}chec de livraison et manque d'outils d'aide \'{a} la d\'{e}cision.

\section{Cadre du stage et objectifs attendus}

Le stage avait un double objectif :
\begin{enumerate}
  \item Concevoir un \textbf{socle CRM op\'{e}rationnel} (gestion des clients, commandes, produits, stock, livraisons et retours).
  \item Int\'{e}grer un \textbf{module d'intelligence artificielle} exploitable en conditions r\'{e}elles pour soutenir la d\'{e}cision m\'{e}tier.
\end{enumerate}

Les livrables valid\'{e}s avec l'organisme d'accueil \'{e}taient :
\begin{itemize}
  \item une application web utilisable avec gestion multi-tenant,
  \item une API backend s\'{e}curis\'{e}e et document\'{e}e,
  \item un microservice ML d\'{e}ployable ind\'{e}pendamment,
  \item une documentation technique et m\'{e}thodologique tra\c{c}able.
\end{itemize}

\section{Organisation de l'\'{e}quipe (4\`{e}me ann\'{e}e)}

Le projet a \'{e}t\'{e} men\'{e} par quatre \'{e}tudiants de 4\`{e}me ann\'{e}e, avec une r\'{e}partition explicite des responsabilit\'{e}s pour garantir une production coh\'{e}rente :

\begin{table}[H]
\centering
\caption{R\'{e}partition des responsabilit\'{e}s principales}
\label{tab:team_roles}
\begin{tabularx}{\textwidth}{l l X}
\toprule
\textbf{\'{E}tudiant} & \textbf{Profil} & \textbf{P\'{e}rim\`{e}tre principal} \\
\midrule
Rezaiguia Soltane Tadj Eddine & Data Science & Pr\'{e}paration des donn\'{e}es, \textit{feature engineering}, analyse des m\'{e}triques et interpr\'{e}tation des r\'{e}sultats. \\
Bellatreche Mohamed Amine & Data Science & Exp\'{e}rimentations comparatives, \'{e}valuation des mod\`{e}les, structuration des jeux d'entra\^{i}nement/test. \\
Khelifi Ayyoub & Computer Science (Cybers\'{e}curit\'{e}) & Backend PHP, architecture applicative, audit de s\'{e}curit\'{e} et impl\'{e}mentation des correctifs. \\
Brahim Soheib & Artificial Intelligence & Int\'{e}gration IA/ML, orchestration FastAPI, endpoints d'insights et flux d'interconnexion avec le CRM. \\
\bottomrule
\end{tabularx}
\end{table}

\section{M\'{e}thode de travail}

L'organisation du stage a suivi une logique incr\'{e}mentale :
\begin{itemize}
  \item d\'{e}composition du projet en lots fonctionnels (CRM, IA, s\'{e}curit\'{e}, int\'{e}gration),
  \item validation progressive par versions,
  \item revues crois\'{e}es entre membres pour limiter les r\'{e}gressions,
  \item consolidation finale par compilation documentaire et v\'{e}rification technique.
\end{itemize}

Cette m\'{e}thode a permis de conserver un fil directeur commun entre les parties m\'{e}tier, techniques et s\'{e}curit\'{e}, sans perdre la sp\'{e}cificit\'{e} de chaque domaine.


% ══════════════════════════════════════════════════════════════
% CHAPITRE 1 : CONTEXTE ET PROBLEMATIQUE
% ══════════════════════════════════════════════════════════════
\chapter{Probl\'{e}matique : Contexte, Objectifs et Contraintes}

\section{Le e-commerce en Alg\'{e}rie}

L'Alg\'{e}rie poss\`{e}de un march\'{e} e-commerce en pleine expansion, port\'{e} par une population jeune et connect\'{e}e. Cependant, plusieurs sp\'{e}cificit\'{e}s distinguent ce march\'{e} :

\begin{itemize}
  \item \textbf{Le mod\`{e}le COD dominant} : En l'absence de solutions de paiement \'{e}lectronique largement adopt\'{e}es, plus de 90\% des transactions se font en Cash-on-Delivery.
  \item \textbf{La g\'{e}ographie logistique} : Le territoire est divis\'{e} en 58 wilayas, regroup\'{e}es en 3 zones d'exp\'{e}dition avec des co\^{u}ts et d\'{e}lais variables.
  \item \textbf{La monnaie locale} : Toutes les transactions s'effectuent en Dinar Alg\'{e}rien (DZD).
\end{itemize}

\section{Probl\'{e}matique du COD}

Le mod\`{e}le COD pr\'{e}sente des d\'{e}fis sp\'{e}cifiques :

\begin{enumerate}
  \item \textbf{Taux d'\'{e}chec \'{e}lev\'{e}} : 30 \`{a} 50\% des commandes COD ne sont jamais livr\'{e}es avec succ\`{e}s (refus, absence, adresse incorrecte).
  \item \textbf{Co\^{u}ts de retour} : Chaque commande \'{e}chou\'{e}e engendre des frais de transport aller-retour sans recette.
  \item \textbf{Gestion manuelle} : Les confirmations t\'{e}l\'{e}phoniques et le suivi sont souvent manuels et chronophages.
  \item \textbf{D\'{e}cisions non inform\'{e}es} : Les gestionnaires manquent d'outils analytiques pour prioriser les commandes ou anticiper les tendances.
\end{enumerate}

\section{Objectifs du projet}

Les objectifs de ce projet sont :

\begin{enumerate}
  \item \textbf{Concevoir un CRM fonctionnel} : Gestion compl\`{e}te des clients, commandes, produits, stocks, livraisons et retours, avec une architecture multi-tenant.
  \item \textbf{Concevoir rigoureusement la base de donn\'{e}es} : Sch\'{e}ma normalis\'{e}, audit trails, int\'{e}grit\'{e} r\'{e}f\'{e}rentielle, support multilingue.
  \item \textbf{Mettre en place un tableau de bord analytique} : Visualisation des KPIs, tendances, et m\'{e}triques en temps r\'{e}el.
  \item \textbf{D\'{e}velopper et comparer des mod\`{e}les de ML} : Au moins deux mod\`{e}les pour la pr\'{e}diction, avec \'{e}valuation rigoureuse.
  \item \textbf{Interpr\'{e}ter les r\'{e}sultats} : Fournir des recommandations actionnables pour l'aide \`{a} la d\'{e}cision.
\end{enumerate}

\section{Contraintes du projet}

Le projet a \'{e}t\'{e} men\'{e} sous plusieurs contraintes structurantes :

\begin{enumerate}
  \item \textbf{Contraintes de donn\'{e}es} : absence de dataset COD alg\'{e}rien complet et exploitable au d\'{e}marrage.
  \item \textbf{Contraintes m\'{e}tier} : n\'{e}cessit\'{e} d'int\'{e}grer les sp\'{e}cificit\'{e}s locales (58 wilayas, DZD, calendrier alg\'{e}rien et islamique).
  \item \textbf{Contraintes syst\`{e}me} : interop\'{e}rabilit\'{e} entre un backend PHP, un frontend Next.js et un microservice Python.
  \item \textbf{Contraintes de s\'{e}curit\'{e}} : protection des donn\'{e}es clients, isolation multi-tenant et durcissement des endpoints expos\'{e}s.
  \item \textbf{Contraintes temporelles} : r\'{e}alisation dans la fen\^{e}tre du stage, avec des jalons de validation acad\'{e}mique.
\end{enumerate}

\section{Crit\`{e}res de r\'{e}ussite retenus}

\begin{table}[H]
\centering
\caption{Crit\`{e}res de r\'{e}ussite m\'{e}tier et technique}
\label{tab:success_criteria}
\begin{tabularx}{\textwidth}{l X X}
\toprule
\textbf{Axe} & \textbf{Crit\`{e}re} & \textbf{Indicateur de validation} \\
\midrule
CRM op\'{e}rationnel & Flux commande COD couvert de bout en bout & Modules commandes/livraison/retours utilisables et reli\'{e}s \\
Architecture & D\'{e}couplage front-back-ML stable & APIs REST fonctionnelles et d\'{e}ploiement ind\'{e}pendant \\
IA & Pr\'{e}dictions exploitables par les agents & Scores de risque, segments clients, pr\'{e}visions disponibles en interface \\
S\'{e}curit\'{e} & Exposition contr\^{o}l\'{e}e des services & Authentification, RBAC, tenant scoping, protections anti-injection \\
Tra\c{c}abilit\'{e} & Possibilit\'{e} de mise \`{a} jour des mod\`{e}les & R\'{e}entra\^{i}nement via base de donn\'{e}es ou CSV avec m\'{e}triques persist\'{e}es \\
\bottomrule
\end{tabularx}
\end{table}


% ══════════════════════════════════════════════════════════════
% CHAPITRE 2 : ETAT DE L'ART
% ══════════════════════════════════════════════════════════════
\chapter{\'{E}tat de l'Art}

\section{CRM intelligent et aide \`{a} la d\'{e}cision}

Un CRM intelligent d\'{e}passe la simple gestion op\'{e}rationnelle : il transforme les historiques de commandes en d\'{e}cisions concr\`{e}tes (priorisation des validations, ciblage commercial et anticipation de la demande). Dans ce projet, l'IA est int\'{e}gr\'{e}e comme couche d'aide \`{a} la d\'{e}cision, sans remplacer les r\`{e}gles m\'{e}tier du CRM.

\section{Techniques ML retenues}

\subsection{Classification du risque de livraison}

La pr\'{e}diction de risque est trait\'{e}e comme une classification binaire (succ\`{e}s \textit{vs} \'{e}chec de livraison). Les trois algorithmes CatBoost \cite{prokhorenkova2018catboost}, LightGBM \cite{ke2017lightgbm} et XGBoost \cite{chen2016xgboost} ont \'{e}t\'{e} retenus car ils restent les plus robustes sur donn\'{e}es tabulaires \cite{grinsztajn2022tree}. Le syst\`{e}me combine leurs scores via un ensemble pond\'{e}r\'{e} afin de limiter la variance.

Le d\'{e}s\'{e}quilibre de classes a \'{e}t\'{e} trait\'{e} par ADASYN \cite{he2008adasyn}, plus adapt\'{e} que SMOTE \cite{chawla2002smote} pour renforcer les zones difficiles de la fronti\`{e}re de d\'{e}cision.

\subsection{Segmentation client}

La segmentation repose sur les variables RFM (Recency, Frequency, Monetary). L'impl\'{e}mentation suit une strat\'{e}gie hybride : HDBSCAN \cite{mcinnes2017hdbscan} est tent\'{e} en priorit\'{e} pour les gros volumes, puis KMeans est appliqu\'{e} lorsque la qualit\'{e} de cluster n'est pas suffisante (silhouette faible ou bruit excessif). Cette logique rend le module exploitable pour des tailles d'entreprise diff\'{e}rentes.

\subsection{Pr\'{e}vision de la demande}

Le benchmark a compar\'{e} des approches classiques et modernes (baseline interne, Prophet \cite{taylor2018forecasting}, Chronos \cite{ansari2024chronos}, LightGBM). LightGBM a \'{e}t\'{e} retenu en production gr\^{a}ce \`{a} son compromis entre performance, temps de r\'{e}entra\^{i}nement et facilit\'{e} d'int\'{e}gration de covariables calendaires alg\'{e}riennes (jours f\'{e}ri\'{e}s nationaux et \'{e}v\'{e}nements islamiques).

\section{M\'{e}triques d'\'{e}valuation}

Les mod\`{e}les de risque sont suivis avec AUC-ROC, pr\'{e}cision, rappel et F1-Score. Les mod\`{e}les de pr\'{e}vision sont \'{e}valu\'{e}s via MAE et RMSE. Ce choix de m\'{e}triques est align\'{e} sur l'objectif m\'{e}tier : r\'{e}duire les \'{e}checs COD tout en maintenant une planification logistique stable.


% ══════════════════════════════════════════════════════════════
% CHAPITRE 3 : CONCEPTION ET ARCHITECTURE
% ══════════════════════════════════════════════════════════════
\chapter{Conception et Architecture}

\section{Architecture globale du syst\`{e}me}

Le syst\`{e}me adopt\'{e} suit une architecture en microservices, compos\'{e}e de trois composants principaux :

\begin{figure}[H]
\centering
\includegraphics[width=\textwidth]{architecture_globale.drawio.png}
\caption{Architecture globale du syst\`{e}me}
\label{fig:architecture}
\end{figure}

Le syst\`{e}me repose sur une communication REST/JSON entre les composants. Le backend PHP (Clean Architecture, pattern Controller-Service-Repository) orchestre les \'{e}changes avec la base MySQL 8.0 (UTF8MB4 pour le multilingue arabe/fran\c{c}ais) et d\'{e}l\`{e}gue les traitements analytiques au microservice ML Python (FastAPI). Ce dernier int\`{e}gre un mod\`{e}le LightGBM avec covariables du calendrier alg\'{e}rien (islamique + jours f\'{e}ri\'{e}s nationaux) pour la pr\'{e}vision de la demande, ainsi que l'API Google Gemini (\texttt{gemini-2.5-flash}) pour la g\'{e}n\'{e}ration d'insights en langage naturel.

\section{Conception de la base de donn\'{e}es}

La base de donn\'{e}es MySQL comprend \textbf{12 tables} r\'{e}parties en \textbf{6 fichiers de migration}, avec une architecture multi-tenant o\`{u} chaque table op\'{e}rationnelle est isol\'{e}e par \texttt{store\_id}.

\subsection{Sch\'{e}ma relationnel}

\begin{figure}[H]
\centering
\includegraphics[width=\textwidth]{schema_base_donnees.drawio.png}
\caption{Sch\'{e}ma relationnel de la base de donn\'{e}es}
\label{fig:schema}
\end{figure}

\subsection{Tables principales}

Les tables couvrent quatre noyaux :
\begin{itemize}
  \item gouvernance multi-tenant (\texttt{stores}, \texttt{users}, \texttt{super\_admins}),
  \item r\'{e}f\'{e}rentiels m\'{e}tier (\texttt{wilayas}, \texttt{products}, \texttt{inventory}),
  \item cycle transactionnel (\texttt{orders}, \texttt{order\_items}, \texttt{deliveries}, \texttt{returns}),
  \item tra\c{c}abilit\'{e} (\texttt{order\_status\_history}, \texttt{inventory\_movements}).
\end{itemize}

\subsection{Cycle de vie d'une commande COD}

La table \texttt{orders} utilise un champ \texttt{status} de type ENUM avec 9 \'{e}tats possibles refl\'{e}tant le cycle de vie complet d'une commande COD :

\begin{figure}[H]
\centering
\includegraphics[width=\textwidth]{cycle_vie_commande.drawio.png}
\caption{Cycle de vie d'une commande COD avec int\'{e}gration IA}
\label{fig:order_states}
\end{figure}

\section{Architecture du module ML}

Le module ML est un microservice Python autonome communiquant avec le backend PHP via API REST.

\subsection{Int\'{e}gration IA dans le syst\`{e}me}

La figure \ref{fig:ai_integration} pr\'{e}sente les trois flux d'analyse IA int\'{e}gr\'{e}s dans le CRM, ainsi que le composant LLM pour la g\'{e}n\'{e}ration d'insights.

\begin{figure}[H]
\centering
\includegraphics[width=\textwidth]{ai_integration.drawio.png}
\caption{Flux d'int\'{e}gration IA dans le CRM}
\label{fig:ai_integration}
\end{figure}

\subsection{Structure du microservice}

Le microservice est structur\'{e} en trois couches :
\begin{itemize}
  \item \textbf{Routes} : exposition des endpoints de pr\'{e}diction, segmentation, pr\'{e}vision, insights et r\'{e}entra\^{i}nement,
  \item \textbf{Services} : orchestration des appels, validation des entr\'{e}es et interfa\c{c}age LLM,
  \item \textbf{Mod\`{e}les} : pipeline d'entra\^{i}nement et artefacts persist\'{e}s (\texttt{trained\_models}).
\end{itemize}

\subsection{Endpoints de l'API ML}

\begin{table}[H]
\centering
\caption{Endpoints op\'{e}rationnels du module ML}
\label{tab:endpoints}
\begin{tabularx}{\textwidth}{l l X}
\toprule
\textbf{M\'{e}thode} & \textbf{Endpoint} & \textbf{Description} \\
\midrule
POST & \texttt{/api/predict/order-risk} & Pr\'{e}diction du risque pour une commande \\
GET & \texttt{/api/segment/customers} & Segmentation des clients \\
GET & \texttt{/api/forecast/demand} & Pr\'{e}vision de la demande \\
GET & \texttt{/api/insights/summary} & Insights g\'{e}n\'{e}r\'{e}s par LLM \\
GET & \texttt{/api/insights/order-explanation} & Explication textuelle d'un score de risque \\
POST & \texttt{/api/retrain/upload-and-train} & R\'{e}entra\^{i}ner via CSV \\
POST & \texttt{/api/retrain/from-database} & R\'{e}entra\^{i}ner depuis la base de donn\'{e}es CRM \\
GET & \texttt{/api/retrain/metrics} & M\'{e}triques d'\'{e}valuation \\
\bottomrule
\end{tabularx}
\end{table}

\subsection{Int\'{e}gration avec le backend PHP}

Le backend PHP encapsule les appels ML dans un service applicatif unique (\texttt{AIInsightsService}), ce qui centralise la gestion des erreurs, les d\'{e}lais de r\'{e}ponse et l'authentification entre services. Le r\'{e}entra\^{i}nement peut ainsi \^{e}tre d\'{e}clench\'{e} depuis l'interface CRM puis trac\'{e} dans les m\'{e}triques de sortie.


% ══════════════════════════════════════════════════════════════
% CHAPITRE 4 : REALISATION
% ══════════════════════════════════════════════════════════════
\chapter{R\'{e}alisation}

\section{Technologies utilis\'{e}es}

\begin{table}[H]
\centering
\caption{Stack technologique du projet}
\label{tab:stack}
\begin{tabularx}{\textwidth}{l l X}
\toprule
\textbf{Composant} & \textbf{Technologie} & \textbf{Version / D\'{e}tails} \\
\midrule
Frontend & Next.js (React) & 14+ avec TypeScript \\
Backend & PHP & 8.2+ (Clean Architecture) \\
Base de donn\'{e}es & MySQL & 8.0+ (UTF8MB4) \\
Module ML & Python / FastAPI & 3.11+ / FastAPI + Uvicorn \\
Moteurs ML & CatBoost, LightGBM, XGBoost, HDBSCAN/KMeans & Apprentissage supervis\'{e} et segmentation \\
LLM & Google Gemini API & Restitution textuelle d'insights \\
Conteneurisation & Docker & Python 3.11-slim \\
\bottomrule
\end{tabularx}
\end{table}

\section{Le CRM : Modules fonctionnels}

Le CRM couvre l'ensemble du cycle COD : authentification et r\^{o}les, catalogue et stock, commandes/livraisons/retours, analytics et administration. Cette couverture op\'{e}rationnelle est coupl\'{e}e au scoping multi-tenant (\texttt{store\_id}) et \`{a} la journalisation des \'{e}tats sensibles.

\section{Jeu de donn\'{e}es : Olist Brazilian E-Commerce}

En l'absence de donn\'{e}es alg\'{e}riennes compl\`{e}tes au d\'{e}marrage, le projet a utilis\'{e} Olist \cite{olist2018} comme socle de benchmark (plus de 100 000 commandes r\'{e}elles, 2016--2018).

\subsection{Transformation et mapping}

Le pipeline applique un mapping vers le contexte local (wilayas, devise DZD, statuts CRM harmonis\'{e}s, cat\'{e}gories normalis\'{e}es). Apr\`{e}s nettoyage, les commandes finalis\'{e}es sont utilis\'{e}es pour l'entra\^{i}nement du mod\`{e}le de risque, afin d'\'{e}viter les \'{e}tiquettes ambigu\"{e}s.

\section{R\'{e}entra\^{i}nement des mod\`{e}les}

Le r\'{e}entra\^{i}nement est activable par deux voies compl\'{e}mentaires :
\begin{itemize}
  \item depuis le backend CRM (\texttt{/api/v1/ai/retrain}),
  \item directement via le service ML (\texttt{/api/retrain/from-database} ou \texttt{/api/retrain/upload-and-train}).
\end{itemize}

Le pipeline suit les \'{e}tapes suivantes : extraction/validation des donn\'{e}es, optimisation Optuna, entra\^{i}nement des trois modules (risque, segmentation, pr\'{e}vision), puis persistance des m\'{e}triques dans \texttt{metrics.json}. Cette cha\^{i}ne garantit la tra\c{c}abilit\'{e} et la reproductibilit\'{e} des versions de mod\`{e}les.


% ══════════════════════════════════════════════════════════════
% CHAPITRE 6 : SOLUTION PROPOSEE (IA + RESULTATS)
% ══════════════════════════════════════════════════════════════
\chapter{Solution Propos\'{e}e : IA, R\'{e}sultats et Validation}

Ce chapitre synth\'{e}tise la logique suivie pendant le stage : benchmark initial sur Olist, puis adaptation aux donn\'{e}es CRM afin d'obtenir une solution d\'{e}ployable en contexte op\'{e}rationnel.

\section{Pipeline global et logique d'int\'{e}gration}

La figure \ref{fig:ml_pipeline} montre la cha\^{i}ne compl\`{e}te (pr\'{e}paration, entra\^{i}nement, validation, d\'{e}ploiement), en coh\'{e}rence avec l'architecture applicative (figure \ref{fig:architecture}) et les flux d'int\'{e}gration IA (figure \ref{fig:ai_integration}).

\begin{figure}[H]
\centering
\includegraphics[width=\textwidth]{ml_pipeline.drawio.png}
\caption{Pipeline Machine Learning : de la donn\'{e}e brute au d\'{e}ploiement}
\label{fig:ml_pipeline}
\end{figure}

\section{D\'{e}marche de d\'{e}veloppement : Olist puis CRM}

La phase Olist a servi \`{a} figer les choix algorithmiques et le protocole d'\'{e}valuation. La phase CRM (base interne ou CSV) a ensuite valid\'{e} l'op\'{e}rabilit\'{e} r\'{e}elle : extraction des donn\'{e}es, r\'{e}entra\^{i}nement, publication des m\'{e}triques et exploitation dans l'interface.

\section{R\'{e}sultats des modules IA}

\subsection{Pr\'{e}diction du risque de livraison}

Le module de risque utilise 25 variables m\'{e}tier et un ensemble CatBoost/LightGBM/XGBoost. Le protocole conserve un split stratifi\'{e}, ADASYN pour le d\'{e}s\'{e}quilibre et une optimisation Optuna.

\begin{table}[H]
\centering
\caption{Performances comparatives (benchmark Olist)}
\label{tab:risk_benchmark_compact}
\begin{tabularx}{\textwidth}{l *{3}{>{\centering\arraybackslash}X}}
\toprule
\textbf{Mod\`{e}le} & \textbf{AUC-ROC} & \textbf{Accuracy} & \textbf{F1} \\
\midrule
CatBoost & 0.9957 & 0.9997 & 0.9999 \\
LightGBM & 0.9974 & 0.9997 & 0.9999 \\
XGBoost & 0.9967 & 0.9996 & 0.9998 \\
\textbf{Ensemble} & \textbf{0.9961} & \textbf{0.9997} & \textbf{0.9999} \\
\bottomrule
\end{tabularx}
\end{table}

Sur un r\'{e}entra\^{i}nement CRM (6 042 commandes), l'ensemble atteint AUC 0.7401 et F1 0.8455. Le score est converti en trois niveaux d'action (faible, interm\'{e}diaire, \'{e}lev\'{e}) pour guider la validation COD.

\subsection{Segmentation client (RFM)}

La segmentation RFM applique une logique hybride HDBSCAN/KMeans. Sur la derni\`{e}re ex\'{e}cution CRM, KMeans est retenu avec 4 segments : VIP (3,1\%), Loyal (11,5\%), Regular (54,6\%) et At Risk (30,9\%). Cette sortie est utilis\'{e}e pour les campagnes de fid\'{e}lisation et de r\'{e}activation.

\subsection{Pr\'{e}vision de la demande}

Le forecasting retient LightGBM avec covariables calendaires alg\'{e}riennes (week-end local, jours f\'{e}ri\'{e}s, Ramadan et A\"{i}d). Le benchmark Olist reste favorable \`{a} cette approche.

\begin{table}[H]
\centering
\caption{Benchmark forecasting (phase Olist, horizon 30 jours)}
\label{tab:forecast_benchmark_compact}
\begin{tabularx}{\textwidth}{l *{3}{>{\centering\arraybackslash}X}}
\toprule
\textbf{Mod\`{e}le} & \textbf{MAE (DZD)} & \textbf{RMSE (DZD)} & \textbf{Ecart vs baseline} \\
\midrule
Moyenne mobile (7j) & 363~075 & 466~570 & baseline \\
Chronos-T5-Small & 338~231 & 395~363 & +6,8\% (MAE) \\
Prophet + jours islamiques & 376~292 & 447~736 & $-$3,6\% (MAE) \\
\textbf{LightGBM + covariables} & \textbf{318~741} & \textbf{416~501} & \textbf{+12,2\% (MAE)} \\
\bottomrule
\end{tabularx}
\end{table}

Sur CRM, la derni\`{e}re ex\'{e}cution indique MAE 47 111 DZD et RMSE 59 129 DZD, avec une am\'{e}lioration MAE de 3,1\% par rapport \`{a} la baseline interne.

\section{Int\'{e}gration LLM et restitution utilisateur}

Gemini est utilis\'{e} pour la restitution explicative (\texttt{/api/insights/summary}, \texttt{/api/insights/order-explanation}, \texttt{/api/insights/recommendations}) et non pour le scoring principal. Les r\'{e}sultats sont pr\'{e}sent\'{e}s en arabe, fran\c{c}ais et anglais.


% ══════════════════════════════════════════════════════════════
% CHAPITRE 7 : SECURITE
% ══════════════════════════════════════════════════════════════
\chapter{S\'{e}curit\'{e} Applicative et Protection des Donn\'{e}es}

Le volet cybers\'{e}curit\'{e} formalise les risques identifi\'{e}s et les correctifs appliqu\'{e}s dans le code afin de garantir un d\'{e}ploiement plus robuste.

\section{M\'{e}thodologie d'audit}

\begin{enumerate}
  \item Revue statique du backend PHP, du frontend Next.js et du microservice FastAPI.
  \item Tests dynamiques via Burp Suite (authentification, contr\^{o}les d'acc\`{e}s, manipulation de requ\^{e}tes).
  \item Campagnes d'injection (sqlmap + tests manuels).
  \item V\'{e}rification de la configuration sensible (secrets, debug, CORS, uploads).
\end{enumerate}

\section{Architecture de protection retenue}

\begin{table}[H]
\centering
\caption{M\'{e}canismes de protection par couche}
\label{tab:security_layers_compact}
\begin{tabularx}{\textwidth}{l X}
\toprule
\textbf{Couche} & \textbf{Contr\^{o}les appliqu\'{e}s} \\
\midrule
Backend PHP & JWT, RBAC, tenant scoping par \texttt{store\_id}, middleware CSRF, rate limiting. \\
ML Service FastAPI & V\'{e}rification \texttt{X-API-Key}/Bearer via \texttt{verify\_api\_key} sur tous les routers. \\
Frontend & \texttt{AuthGuard} sur routes prot\'{e}g\'{e}es, gestion des erreurs de rendu. \\
Donn\'{e}es et configuration & Requ\^{e}tes param\'{e}tr\'{e}es, suppression de secrets versionn\'{e}s, templates \texttt{.env.example}. \\
\bottomrule
\end{tabularx}
\end{table}

\section{Synth\`{e}se des corrections int\'{e}gr\'{e}es}

\begin{table}[H]
\centering
\caption{Synth\`{e}se des vuln\'{e}rabilit\'{e}s corrig\'{e}es (V-01 \`{a} V-10)}
\label{tab:security_summary_compact}
\begin{tabularx}{\textwidth}{l X}
\toprule
\textbf{Famille de risque} & \textbf{Correction appliqu\'{e}e dans le projet} \\
\midrule
Exposition de secrets & Nettoyage des cl\'{e}s sensibles dans les fichiers suivis, configuration par variables d'environnement. \\
Acc\`{e}s non autoris\'{e} / IDOR & Contr\^{o}les d'acc\`{e}s renforc\'{e}s (RBAC + scoping tenant) et garde frontend. \\
Injection SQL & Param\'{e}trisation syst\'{e}matique des requ\^{e}tes, y compris pagination. \\
Attaques d'\'{e}tat (CSRF) & Middleware CSRF + en-t\^{e}te applicatif c\^{o}t\'{e} client. \\
Abus d'API / flooding & Rate limiting sur groupes de routes sensibles. \\
Upload malveillant & Validation type/taille/nom et blocage de path traversal. \\
\bottomrule
\end{tabularx}
\end{table}

\section{Pr\'{e}conisations de d\'{e}ploiement}

\begin{itemize}
  \item imposer HTTPS sur toutes les couches expos\'{e}es,
  \item activer la rotation p\'{e}riodique des secrets (JWT, ML API key, LLM key),
  \item journaliser les rejets de s\'{e}curit\'{e} (auth, CSRF, rate-limit),
  \item planifier un pentest externe avant d\'{e}ploiement \`{a} grande \'{e}chelle.
\end{itemize}


% ══════════════════════════════════════════════════════════════
% CHAPITRE 8 : REVUE FINALE
% ══════════════════════════════════════════════════════════════
\section{Revue finale de coh\'{e}rence}

Une v\'{e}rification crois\'{e}e a \'{e}t\'{e} men\'{e}e entre sch\'{e}mas, routes et impl\'{e}mentation effective. Les points suivants sont align\'{e}s :
\begin{itemize}
  \item architecture Frontend-Backend-ML conforme aux sch\'{e}mas \texttt{architecture\_globale} et \texttt{ai\_integration},
  \item mod\`{e}le de donn\'{e}es conforme aux migrations (12 tables m\'{e}tier et isolation multi-tenant),
  \item cycle de vie des commandes coh\'{e}rent avec le \texttt{ENUM status},
  \item endpoints de r\'{e}entra\^{i}nement et d'insights conformes au code actuel.
\end{itemize}

Les ajustements documentaires finaux ont aussi \'{e}t\'{e} valid\'{e}s, notamment la correction de la page de garde et l'harmonisation du vocabulaire utilis\'{e} dans la partie s\'{e}curit\'{e}.


% ══════════════════════════════════════════════════════════════
% CONCLUSION GENERALE
% ══════════════════════════════════════════════════════════════
\chapter*{Conclusion G\'{e}n\'{e}rale}
\addcontentsline{toc}{chapter}{Conclusion G\'{e}n\'{e}rale}

Ce projet a abouti \`{a} un CRM intelligent exploitable pour le contexte COD alg\'{e}rien, avec une cha\^{i}ne compl\`{e}te de la donn\'{e}e brute jusqu'\`{a} la d\'{e}cision op\'{e}rationnelle.

\textbf{Synth\`{e}se des apports :}
\begin{itemize}
  \item architecture multi-tenant coh\'{e}rente (frontend, backend, microservice ML),
  \item d\'{e}marche m\'{e}thodologique r\'{e}aliste (benchmark Olist puis adaptation CRM),
  \item trois modules IA int\'{e}gr\'{e}s (risque, segmentation, pr\'{e}vision) avec indicateurs mesurables,
  \item s\'{e}curisation applicative renforc\'{e}e (authentification, contr\^{o}les d'acc\`{e}s, protections anti-abus).
\end{itemize}

\textbf{Perspectives :}
\begin{itemize}
  \item multiplier les r\'{e}entra\^{i}nements sur donn\'{e}es r\'{e}elles pour stabiliser les performances en production,
  \item enrichir les covariables m\'{e}tier (saisonnalit\'{e} commerciale, campagnes marketing, effets logistiques),
  \item industrialiser l'exploitation (CI/CD, supervision avanc\'{e}e, tests de s\'{e}curit\'{e} externes).
\end{itemize}


% ══════════════════════════════════════════════════════════════
% BIBLIOGRAPHIE
% ══════════════════════════════════════════════════════════════
\begin{thebibliography}{99}
\addcontentsline{toc}{chapter}{Bibliographie}

\bibitem{chawla2002smote}
N. V. Chawla, K. W. Bowyer, L. O. Hall, and W. P. Kegelmeyer, ``SMOTE: Synthetic Minority Over-sampling Technique,'' \textit{Journal of Artificial Intelligence Research}, vol. 16, pp. 321--357, 2002.

\bibitem{he2008adasyn}
H. He, Y. Bai, E. A. Garcia, and S. Li, ``ADASYN: Adaptive Synthetic Sampling Approach for Imbalanced Learning,'' in \textit{Proceedings of the IEEE International Joint Conference on Neural Networks}, pp. 1322--1328, 2008.

\bibitem{chen2016xgboost}
T. Chen and C. Guestrin, ``XGBoost: A Scalable Tree Boosting System,'' in \textit{Proceedings of the 22nd ACM SIGKDD International Conference on Knowledge Discovery and Data Mining}, pp. 785--794, 2016.

\bibitem{ke2017lightgbm}
G. Ke, Q. Meng, T. Finley, T. Wang, W. Chen, W. Ma, Q. Ye, and T.-Y. Liu, ``LightGBM: A Highly Efficient Gradient Boosting Decision Tree,'' in \textit{Advances in Neural Information Processing Systems 30}, pp. 3146--3154, 2017.

\bibitem{prokhorenkova2018catboost}
L. Prokhorenkova, G. Gusev, A. Vorobev, A. V. Dorogush, and A. Gulin, ``CatBoost: Unbiased Boosting with Categorical Features,'' in \textit{Advances in Neural Information Processing Systems 31}, pp. 6638--6648, 2018.

\bibitem{taylor2018forecasting}
S. J. Taylor and B. Letham, ``Forecasting at Scale,'' \textit{The American Statistician}, vol. 72, no. 1, pp. 37--45, 2018.

\bibitem{mcinnes2017hdbscan}
L. McInnes, J. Healy, and S. Astels, ``hdbscan: Hierarchical density based clustering,'' \textit{Journal of Open Source Software}, vol. 2, no. 11, p. 205, 2017.

\bibitem{grinsztajn2022tree}
L. Grinsztajn, E. Oyallon, and G. Varoquaux, ``Why do tree-based models still outperform deep learning on tabular data?'' in \textit{Advances in Neural Information Processing Systems 35}, 2022.

\bibitem{olist2018}
Olist, ``Brazilian E-Commerce Public Dataset by Olist,'' Kaggle, 2018. [Online]. Available: \url{https://www.kaggle.com/datasets/olistbr/brazilian-ecommerce}

\bibitem{ansari2024chronos}
A. F. Ansari, L. Stella, C. Turkmen, X. Zhang, P. Mercado, H. Shen, O. Shchur, S. S. Rangapuram, S. P. Arango, S. Kapoor, J. Zschiegner, D. C. Maddix, M. W. Mahoney, K. Torkkola, A. G. Wilson, M. Bohlke-Schneider, and Y. Wang, ``Chronos: Learning the Language of Time Series,'' \textit{arXiv preprint arXiv:2403.07815}, 2024.

\end{thebibliography}


% ══════════════════════════════════════════════════════════════
% ANNEXES
% ══════════════════════════════════════════════════════════════
\end{document}
